\documentclass[10pt,a4paper]{article}
\usepackage[margin=2.2cm]{geometry}
\usepackage{fontspec}
\usepackage{kotex}
\setmainfont{Apple SD Gothic Neo}
\setsansfont{Apple SD Gothic Neo}
\usepackage{amsmath,amssymb}
\usepackage{graphicx}
\usepackage{booktabs}
\usepackage{caption}
\usepackage[dvipsnames]{xcolor}
\usepackage[colorlinks=true,linkcolor=MidnightBlue,citecolor=MidnightBlue,
            urlcolor=MidnightBlue]{hyperref}
\usepackage{titlesec}
\titleformat{\section}{\large\bfseries}{\thesection}{0.6em}{}
\titleformat{\subsection}{\normalsize\bfseries}{\thesubsection}{0.5em}{}
\captionsetup{font=small,labelfont=bf}
\setlength{\parskip}{0.35em}
\setlength{\parindent}{0pt}

% 주장 상자 — 모든 주장은 반증 조건을 함께 적는다
\usepackage{mdframed}
\newmdenv[linewidth=0.4pt,linecolor=gray!60,backgroundcolor=gray!5,
          innertopmargin=6pt,innerbottommargin=6pt,skipabove=8pt,skipbelow=8pt]{claimbox}
\newcommand{\claim}[2]{\begin{claimbox}\textbf{주장.} #1\par\smallskip
  \textcolor{gray!70}{\textbf{반증 조건.} #2}\end{claimbox}}

\title{\bfseries 덮음은 값이 아니라 \emph{부호}를 정한다}
\author{월드모델 실험실 · 노트 660--661}
\date{2026-08-06}
\begin{document}\maketitle

\section{한 줄}

같은 축, 같은 도메인, 같은 열이다. 덮음이 $0.43$ 일 때 그 도메인은 $-0.0227$ 을
잃고 $0.82$ 일 때 $+0.0058$ 을 번다. \textbf{덮음이 낮으면 값이 작아지는 것이
아니라 부호가 뒤집힌다.}

\section{노트 548 이 세운 것}

노트 548 이 ``덮음은 값이 아니다''를 쟀다 --- 자료를 더 모을지는 덮음이 아니라
그 축이 라벨과 붙는가로 정하고, 순위 상관이 값$\leftrightarrow$상관 $+0.786$ 대
값$\leftrightarrow$덮음 $+0.451$ 이었다. 그 결론은 \emph{수집 우선순위}에 대한
것이었고 맞다.

이 노트는 다른 것을 말한다 --- \textbf{덮음이 낮은 축을 넣으면 무슨 일이
일어나나.}

\section{쟀다}

지역 방문자 축 하나(\verb|vis_out|)를 팝업과 시장팝업이 나눠 쓰게 하고, 시장팝업
덮음만 바꿔 같은 네 팔을 두 번 돌렸다. 기준선과 팝업만 팔이 세 번 다 정확히
재현됐다($0.4689$ · $+0.0008$ · 팝업 $+0.009$ · 8/12).

\begin{center}\small
\begin{tabular}{lrrr}\toprule
팔 & 판 & 팝업 & 시장팝업 \\ \midrule
팝업만 & $+0.0008$ (8/12) & $+0.009$ & $+0.002$ \\
팝업$+$시장 (덮음 $0.43$) & $-0.0005$ (5/12) & $+0.0201$ & $\mathbf{-0.0227}$ \\
팝업$+$시장 (덮음 $\mathbf{0.82}$) & $+0.0002$ (9/12) & $+0.0064$ & $\mathbf{+0.0058}$ \\
위약 (덮음 $0.82$) & $-0.0040$ (0/12) & $-0.1297$ & $-0.0399$ \\ \bottomrule
\end{tabular}\end{center}

\section{위약이 기제를 보여 준다}

덮음 $0.43$ 에서는 진짜와 위약의 \textbf{도메인 부호가 정확히 반대}였다 ---
시장팝업 진짜 $-0.0227$ 대 위약 $+0.0151$. \textbf{값을 섞으면 나아진다.} 값이
없어져도 마스크 무늬는 그대로이므로, 남은 것(결측 표시자)이 유용했다는 뜻이다.

덮음 $0.82$ 에서는 \textbf{둘 다 음수}다(팝업 $-0.1297$ · 시장팝업 $-0.0399$).
값이 두 도메인에서 다 정보다.

\claim{덮음이 낮은 도메인에 열을 더하면 결측 무리가 정보 대신 표시자로 쓰이고,
그것이 부호를 뒤집는다.}{덮음을 올렸는데 부호가 그대로면 이 설명이 틀렸다.}

덮음을 $0.430 \to 0.823$ 으로 올려 부호가 뒤집혔으므로 이 주장은 살았다.
올린 방법은 지리다 --- 시장 레코드는 기사에서 만들어 도로명주소가 없지만
\verb|neighborhood| 에 \textbf{동 이름}이 있다(성수동 · 여의도 · 잠실 · 한남동).
동$\to$구 표 130항목을 손으로 적었고(라벨을 안 본다) 매칭 표본 10건을 눈으로
확인해 10/10 정확했다.

\section{신호 몫이 문턱에 닿았다}

판 $+0.0002$ 는 문턱 $0.0045$ 에 한참 못 미친다 --- \textbf{채택 안 한다}.
그런데 \textbf{신호 몫}(진짜 $-$ 위약)은 $+0.0002-(-0.0040) = \mathbf{+0.0042}$ 이고
문턱에 거의 닿는다. 이 판에서 상태 축이 낸 가장 큰 값이다(노트 650 의 혼잡도
$+0.0021$ 의 두 배).

\section{덤 --- 이웃 효과 여덟 번째}

팝업이 $+0.0201$(덮음 $0.43$ 시절)에서 $+0.0064$(덮음 $0.82$)로 \emph{내려갔다}.
노트 660 에서 ``열을 나눠 쓰니 팝업 이득이 두 배''라고 적은 것은 \textbf{시장팝업의
결측 표시자가 만든 이득}이었을 가능성이 크다 --- 표시자가 사라지니 팝업도 원래
수준($+0.009$)에 가까워졌다.

이 판이 이웃 효과를 여덟 번 만났다(노트 525 $\cdot$ 551 $\cdot$ 559 $\cdot$
560 $\cdot$ 563 $\cdot$ 596 $\cdot$ 610 $\cdot$ 661). \textbf{한 도메인의 이득을
그 도메인의 자료로 설명하기 전에 이웃을 본다.}

\begin{figure}[h]\centering
\includegraphics[width=\textwidth]{figs/f1.pdf}
\caption{(a) 같은 축 · 같은 열인데 덮음이 부호를 뒤집는다(회색 띠는 판 문턱).
(b) 덮음 $0.43$ --- 진짜와 위약의 부호가 반대다(섞으면 나아진다).
(c) 덮음 $0.82$ --- 둘 다 같은 방향이고 신호 몫이 판 $+0.0042$ 다.}
\end{figure}

\end{document}
